\scanpage{81}\mspage{68}
% archive sheet number: 80
% Paragraph continues from the preceding sheet.
admet sans plus que \uline{dans ce cas} $u_!^W$ est
adjoint à gauche de $u_W^*$ ; et \uline{de plus},
que si $u$ est $W_0$-\deleted{\unsure{homotopique}}\inserted{\unsure{asphérique}} donc
$u^*:M(I'^{\unsure{o|\circ}})\unsure{\text{---}|\to}M(I^o)$ pl. fidèle, qu'il en
soit de même de $u_W^{\unsure{'|1|*}}$ déduit par passage
aux localisés. Or même si on admet
que $u_!^W$ est adj. à gauche de $u_{\unsure{*|\times}}^W$,
faute de disposer d'un morphisme
d'adjonction tout prêt ``en bas'' (dans
$\mathrm{Ho}_W(I'^o)$), déduit du morphisme
d'adjonction dans $M(I'^o)$, il n'est
absolument pas clair que \deleted{\unsure{$u_!^W$}}\unsure{ce}
$u_W^*$ soit pl. fid. i.e. $u_!^W$ un foncteur
de localisation. Le contre-exemple 3.6.
montre bien que, $\unsure{=}$ si on a adjonction
(i.e. si $e$ est un objet initial de $C=C'$],
le foncteur $u^*=\mathrm{id}_C$ a beau être pl.
fid (et $\mathord{=}$ une équivalence de catégories),
le foncteur localisé n'a aucune raison
de l'être : cela prouverait que \unsure{tt} objet
\deleted{catégorie}\inserted{de $\mathrm{Cat}$} ayant un objet initial est
\deleted{un groupoïde} équivalent à $e$ !
